\chapter{Solutions to Exercises}\label{app:solutions}

\section*{Chapter 1}

\paragraph{Solution to Exercise \ref{ex:1-null}.}
\emph{Moments.} $\Pi(i)$ is uniform on $\{1,\dots,N\}$, so
$\E[b_{\Pi(i)}^2] = \frac1N \sum_j b_j^2 = \frac1N$ by normalization. For
$i \ne i'$, the pair $(\Pi(i), \Pi(i'))$ is uniform on ordered pairs
$j \ne j'$, so
\[
  \E[b_{\Pi(i)} b_{\Pi(i')}]
  = \frac{1}{N(N-1)} \sum_{j \ne j'} b_j b_{j'}
  = \frac{\bigl(\sum_j b_j\bigr)^2 - \sum_j b_j^2}{N(N-1)}
  = \frac{0 - 1}{N(N-1)} .
\]
\emph{Asymptotic normality.} $\IC = \sum_i a_i b_{\Pi(i)}$ is a simple
linear rank statistic. Hoeffding's combinatorial CLT gives normality
provided no single term dominates (the Noether condition
$\max_i a_i^2 / \sum_i a_i^2 \to 0$). For rank scores,
$a_i \propto i - \frac{N+1}{2}$ with
$\sum_i (i - \frac{N+1}{2})^2 = \frac{N(N^2-1)}{12}$, so
$\max_i a_i^2 = \frac{(N-1)^2/4}{N(N^2-1)/12} = \frac{3(N-1)}{N(N+1)} \to 0$.
Hence $\IC/\sqrt{\Var(\IC)} = \sqrt{N-1}\,\IC \Rightarrow \gauss(0,1)$.

\paragraph{Solution to Exercise \ref{ex:1-quantile}.}
Standardize the score, $Z = s/\sigma_s$. Joint normality gives
$\E[y \mid Z] = \rho\,\sigma_y Z$, so by the tower property the long leg's
expected return is $\rho\,\sigma_y\,\E[Z \mid Z > z_{0.8}]$ and the short
leg's is $\rho\,\sigma_y\,\E[Z \mid Z < z_{0.2}]$. With
$z_{0.8} = 0.8416$, $\phi(z_{0.8}) = 0.2800$:
$\E[Z \mid Z > z_{0.8}] = \phi(z_{0.8})/0.2 = 1.400$, and by symmetry the
lower tail gives $-1.400$. Hence
$\E[\mathrm{L{-}S}] = 2.80\,\rho\,\sigma_y$. Numerically:
$2.80 \times 0.03 \times 0.02 = 1.68\times10^{-3}$, i.e.\
\emph{$\approx 17$ basis points per holding period}, gross. Against
Chapter~1's cost accounting (a fresh two-sided book turns over fully:
$2\times$ turnover $\times$ 10 bp $= 4$ bp if fully rebuilt daily), the
margin is real but thin --- which is why turnover is a headline number, not
a footnote.

\paragraph{Solution to Exercise \ref{ex:1-bestofN}.}
Core simulation (vectorized; ranks via double \code{argsort}):

\begin{lstlisting}
rng = np.random.default_rng(0)
best_t = []
for rep in range(1000):
    ics = np.empty((20, 250))
    for m in range(20):                    # 20 strategies
        p = rng.standard_normal((250, 100))
        y = rng.standard_normal((250, 100))
        rp = p.argsort(1).argsort(1); ry = y.argsort(1).argsort(1)
        rp = rp - rp.mean(1, keepdims=True); ry = ry - ry.mean(1, keepdims=True)
        ics[m] = (rp * ry).sum(1) / np.sqrt((rp**2).sum(1) * (ry**2).sum(1))
    t = ics.mean(1) / (ics.std(1, ddof=1) / np.sqrt(250))
    best_t.append(t.max())
\end{lstlisting}

Expected findings: each strategy's $t$ is approximately standard normal
under the null, so the best of $M = 20$ has
$\Prob(\max t > 2) = 1 - \Phi(2)^{20} = 1 - 0.97725^{20} \approx 0.37$ ---
simulation lands near $37\%$ --- and the mean best-$t$ sits near the
$\sqrt{2\ln 20} \approx 2.45$ heuristic. The best strategy's $\ICbar$ is
about $1.87 \times 0.1/\sqrt{250} \approx 0.012$ (expected maximum of 20
standard normals times the SE): a plausible-looking ``signal,''
manufactured from nothing, better than one date-per-week of real research
projects. This is Remark~\ref{rem:ratchet} in numbers, and the reason
Chapter~17's registry counts trials.

\paragraph{Solution to Exercise \ref{ex:1-tstat}.}
(a) $T \ge (3 \times 0.12/0.02)^2 = 18^2 = 324$ dates. (b) If the IC series
has autocorrelations $\rho_s = 1 - s/h$ for $s < h$ and $0$ beyond
(the linear-decay overlap profile), then for $T \gg h$,
\[
  \Var(\ICbar) \approx \frac{\sigma_{\IC}^2}{T}
  \Bigl(1 + 2\sum_{s=1}^{h-1}\bigl(1 - \tfrac{s}{h}\bigr)\Bigr)
  = \frac{\sigma_{\IC}^2}{T}\bigl(1 + (h-1)\bigr)
  = \frac{\sigma_{\IC}^2}{T}\, h,
\]
using $\sum_{s=1}^{h-1}(1-s/h) = (h-1)/2$. The variance inflation factor is
exactly $h$: overlapping $5$-day labels make $T$ dates worth $T/5$
independent ones. Corrected requirement: $324 \times 5 = 1{,}620$ dates ---
about $6.5$ years of daily out-of-sample. (c) Any
\code{n\_folds} $\times$ \code{test\_dates\_per\_fold} $\ge 1{,}620$ works
arithmetically, e.g.\ $9 \times 180$; on a $2{,}500$-date panel this leaves
$\approx 880$ dates for the first fold's training window --- the practical
tension (evaluation appetite versus training data) that walk-forward design
always negotiates, and a reason thesis-scale claims favor pooled multi-year
evaluations over quarterly ones.

\paragraph{Solution to Exercise \ref{ex:1-nll}.}
(a) $\partial \NLL/\partial(\sigma^2) =
\frac{1}{2\sigma^2} - \frac{1}{2\sigma^4}\cdot\frac1n\sum_j e_j^2 = 0$
gives $\hat\sigma^2 = \frac1n\sum_j e_j^2$; substituting back yields
$\tfrac12\log\hat\sigma^2 + \tfrac12\log 2\pi + \tfrac12$.
(b) $\Delta\NLL = \tfrac12 \log(\hat\sigma_A^2/\hat\sigma_B^2)$, so
$\hat\sigma_A^2/\hat\sigma_B^2 = e^{2 \times 0.021} = 1.043$: model A's
residual variance is about $4.3\%$ higher. (c) Two points with
heteroskedastic truth. Model A predicts errors $(0.1, 10)$ but knows its
uncertainty, $\sigma = (0.1, 10)$: per-point NLL (dropping constants)
$= [\log 0.1 + \tfrac12] + [\log 10 + \tfrac12] = 1.0$, MSE $= 50$. Model B
predicts better, errors $(0.05, 8)$, but claims $\sigma = 1$ everywhere:
NLL $= \tfrac12(0.0025) + \tfrac12(64) \approx 32$, MSE $= 32$. B wins MSE
$32 < 50$; A wins NLL $1.0 \ll 32$. The likelihood is rewarding
\emph{calibrated uncertainty} --- being right about how wrong you will be
--- which is precisely the capacity Chapter~2 argues regime models must
have, and MSE cannot see.

\paragraph{Solution to Exercise \ref{ex:1-panel}.}
(a) \code{x\_seq} is $(n, 8, 2)$ and \code{x\_snap} $(n, 3)$ for the test
spec, \code{y} a separate $(n,)$ tensor; no feature column can equal the
target because the target is not stored among features at all.
(b) \code{split\_by\_date(0.8)} yields
$\max(\code{train.date}) < \min(\code{test.date})$; assert it directly.
(c) The public surface offers \code{full\_batch}, \code{minibatches}
(within a panel already split), \code{date\_batches},
\code{subset\_dates}, \code{split\_by\_date} --- every subsetting operation
is keyed by date; row-level selection (\code{\_take}, \code{\_batch}) is
private by underscore convention. The honest path is the paved one.
\code{test\_no\_lookahead} pins the invariant of
Definition~\ref{def:pit}: it rebuilds features after perturbing only
\emph{future} prices and asserts every feature at date $t$ is unchanged
while the forward-return target changes --- look-ahead, if ever introduced,
fails a test rather than improving a backtest.

\section*{Chapter 2}

\paragraph{Solution to Exercise \ref{ex:2-kurt}.}
(a) $\kappa - 3 = 3\bigl(\E[V^2] - (\E V)^2\bigr)/(\E V)^2 =
3\Var(V)/(\E V)^2$. (b) Let $q$ be the weight of the volatile state and
$r = \sigma_2^2/\sigma_1^2 = 25$. Then
$\kappa(q) = 3\,\frac{1 + (r^2-1)q}{(1 + (r-1)q)^2}$; setting
$\kappa'(q) = 0$ gives $(r^2-1)(1 + (r-1)q) = 2(r-1)(1 + (r^2-1)q)$, i.e.\
$q^\ast = \frac{1}{r+1} = \frac{1}{26} \approx 3.8\%$, with
$\kappa(q^\ast) = \frac{3(r+1)^2}{4r} \approx 20.3$. Rare volatile regimes
maximize kurtosis because $\kappa - 3$ is the \emph{relative} variance of
the randomized variance $V$: a scarce, extreme value of $V$ maximizes
$\Var(V)$ against a small mean, exactly as a rare large claim dominates an
insurer's loss distribution. (c) The three numbers say: equal-weight
regime mixing underexplains empirical tails ($5.6$ vs.\ $10$--$30$), but
\emph{realistically proportioned} regimes --- turbulence as the rare state
--- put the same machinery in range ($\approx 20$). Regime switching is
thus a sufficient explanation for most of empirical kurtosis \emph{only if}
the volatile state is genuinely rare and distinct, which is also what the
persistence estimates of Proposition~\ref{prop:acf} deliver
(long calm spells, occasional storms); residual excess beyond that
suggests within-regime non-normality or more than two states.

\paragraph{Solution to Exercise \ref{ex:2-acf}.}
(a) With $a = b$ the stationary law is $(\tfrac12, \tfrac12)$; write
$v \in \{\sigma_1^2, \sigma_2^2\}$. From the proof:
$\Cov(r_t^2, r_{t+s}^2) = \Cov(v_t, v_{t+s}) = \Var(v)\lambda^s$ and
$\Var(r_t^2) = 3\E[v^2] - (\E v)^2$, so
$c = \Var(v)/\bigl(3\E[v^2] - (\E v)^2\bigr)$. For
$\sigma = (0.6, 2.2)$: $v \in \{0.36, 4.84\}$, $\E v = 2.6$,
$\E v^2 = 11.778$, $\Var(v) = 5.018$, so
$c = 5.018/(35.333 - 6.76) = 5.018/28.573 = 0.176$; with
$\lambda = 2(0.98) - 1 = 0.96$, the lag-10 autocorrelation is
$0.176 \times 0.96^{10} = 0.117$. (b) Duration in state 1 is the number of
steps until the first departure: $\Prob(D = n) = a^{n-1}(1-a)$, geometric,
$\E D = 1/(1-a)$. (c) GARCH(1,1) persistence $\alpha + \beta$ absorbs
$\lambda$: the fitted volatility process must reproduce an ACF decaying
like $0.96^s$, and $\alpha+\beta$ \emph{is} the ACF decay rate of a
GARCH's squared returns. This is the classic Lamoureux--Lastrapes point:
ignored regime switching masquerades as near-integrated GARCH
($\alpha + \beta \approx 1$), so extreme persistence estimates can be an
artifact of unmodeled states.

\paragraph{Solution to Exercise \ref{ex:2-pool}.}
(a) $\beta_{\mathrm{pool}} = \E[\beta_S] = p\beta_1 + (1-p)\beta_2$
(independence of $x$ and $S$); since
$\Var(y) = \E[\beta_S^2] + \sigma_\varepsilon^2$ (unit-variance $x$),
$R^2_{\mathrm{pool}} = (\E\beta_S)^2 / (\E[\beta_S^2] +
\sigma_\varepsilon^2)$. (b) Pooled residual:
$y - \beta_{\mathrm{pool}}x = (\beta_S - \E\beta_S)x + \varepsilon$, whose
variance is $\Var(\beta_S)\,\E[x^2] + \sigma_\varepsilon^2$ --- exceeding
the oracle's $\sigma_\varepsilon^2$ by $\Var(\beta_S)\E[x^2]$. (c) Both
results are the law of total variance with a latent index: unexplained
dispersion of a latent \emph{parameter} (variance $\sigma_S^2$ in
Proposition~\ref{prop:kurtosis}, slope $\beta_S$ here) reappears as
inflation of an unconditional moment (kurtosis there, residual variance
here). ``Variance of the variance'' produces fat tails; ``variance of the
slope'' produces phantom noise --- one phenomenon, two disguises, and the
mixture model is the device that gives the latent index somewhere to live
instead.

\paragraph{Solution to Exercise \ref{ex:2-modes}.}
Center and scale: $\mu_{1,2} = \pm\delta$, $\sigma = 1$. The density is
proportional to
$e^{-(x-\delta)^2/2} + e^{-(x+\delta)^2/2}
= 2e^{-(x^2+\delta^2)/2}\cosh(\delta x)$. Setting the log-derivative to
zero: $x = \delta \tanh(\delta x)$. $x = 0$ is always critical. The map
$x \mapsto \delta\tanh(\delta x)$ has slope $\delta^2$ at the origin and is
strictly concave on $x > 0$; nonzero fixed points therefore exist iff
$\delta^2 > 1$. Checking curvature, $x = 0$ is the unique maximum when
$\delta \le 1$ and becomes a local \emph{minimum} flanked by two symmetric
maxima when $\delta > 1$. In original units: bimodal iff
$|\mu_1 - \mu_2| > 2\sigma$. Conclusion: components must be separated by
more than two standard deviations before the mixture \emph{looks} like a
mixture; regime structure at realistic separations is invisible to the
eye, so recovery must be graded against planted latent states
(Chapter~14), not histograms.

\paragraph{Solution to Exercise \ref{ex:2-hands}.}
Sketch of expected results (exact values vary with seed):
(a) sample kurtosis of the pooled series lands near the
$3\,\E[V^2]/(\E V)^2 = 5.56$ prediction for equal-weight
$\sigma = (0.5, 2.5)$ (the Markov chain with $a = b$ spends half its time
in each state); the log-ACF of $r_t^2$ is linear in lag with slope
$\log(0.96) \approx -0.041$, i.e.\ $\approx 4\%$ decay per day.
(b) Pooled OLS slopes on the snapshot features are statistically
indistinguishable from zero while within-regime regressions recover
$\pm\beta$ with $R^2 = \beta^2/(\beta^2 + \sigma_\varepsilon^2)$; with
\code{beta\_scale}${}=2$, \code{noise\_std}${}=0.4$ the oracle $R^2$
exceeds $0.9$ --- the full Proposition~\ref{prop:signflip} contrast.
(c) With \code{regime\_process="iid"} the squared-return ACF collapses to
zero at all positive lags while the kurtosis is unchanged --- mixing makes
tails, persistence makes memory, exactly the division of
Sections~\ref{sec:mixmath}.

\section*{Chapter 3}

\paragraph{Solution to Exercise \ref{ex:3-argmax}.}
\begin{lstlisting}
z = torch.tensor([0.2, 1.0, -0.5], requires_grad=True)
a, b = torch.tensor(3.0), torch.tensor(-3.0)
pred = torch.where(torch.softmax(z, 0).argmax() == 1, a, b)
pred.pow(2).backward()
print(z.grad)   # -> None (no graph reaches z), or zeros if z re-enters
\end{lstlisting}
\code{argmax} returns an integer with no gradient function, and
\code{where} on that condition selects between tensors that do not depend
on $z$; the loss is a piecewise-constant function of $z$, whose derivative
is zero wherever it exists. Replacing the selection with
$\hat y = \sum_k \pi_k \mu_k$ (\code{(torch.softmax(z,0) *
torch.stack([b, a, b])).sum()}) routes real gradient into every logit:
$\partial \hat y/\partial z_k = \pi_k(\mu_k - \hat y)$, alive whenever the
experts disagree. The package's \code{test\_gate\_receives\_gradient}
asserts the second fact; its argmax companion asserts the first ---
pinning both the defect and the repair.

\paragraph{Solution to Exercise \ref{ex:3-params}.}
(a) Each LSTM gate $g \in \{i, f, o, \tilde c\}$ computes
$g = \phi(W_g x_t + U_g h_{t-1} + b_g)$ with $W_g \in \R^{n \times d}$,
$U_g \in \R^{n \times n}$, $b_g \in \R^n$: four gates give
$4\bigl(n(n+d) + n\bigr)$. The GRU has three such blocks ($r$, $z$,
candidate): $3\bigl(n(n+d)+n\bigr)$. (b) Prototype, $d = 5$: layer 1,
$4(256 \cdot 261 + 256) \approx 268$K; layers 2--8 have input dimension
$256$: $4(256 \cdot 512 + 256) \approx 525$K each, $\times 7 \approx
3.68$M; total $\approx 3.95$M. Corrected, GRU $n = 32$, $d = 2$:
$3(32 \cdot 34 + 32) = 3{,}360$. Three orders of magnitude apart.
(c) (i) All rows: $1.25$M samples $\to$ budget $\approx 125$K parameters.
(ii) $N_{\text{eff}} = 20$: $50$K effective samples $\to$ budget
$\approx 5$K. The prototype exceeds even the na\"ive budget $30\times$ and
the discounted one $\approx 800\times$; the corrected encoder fits
\emph{inside} the pessimistic budget. The point of the exercise is that
this verdict required no training run --- it is design-time arithmetic.

\paragraph{Solution to Exercise \ref{ex:3-bn}.}
Build eight samples for date 1 ($\gauss(0,1)$ features) and eight for date
2 with a common $+5$ shift. \code{BatchNorm1d} in train mode, forward on
date 1 alone versus on the concatenated batch: date-1 outputs differ
(the batch mean absorbs date 2's shift --- with the shift included, the
batch mean rises by $2.5$, so date-1 samples are pushed down), a
back-in-time flow of date-2 information: consequence (iii), leakage, via
consequence (ii), batch-composition sensitivity. Switching the layer to
\code{eval()} changes date-1 outputs yet again (running statistics replace
batch statistics): consequence (i), train/eval asymmetry.
\code{LayerNorm} normalizes each sample by \emph{its own} feature moments,
so date-1 outputs are bit-identical with or without date 2 present ---
\code{torch.equal} returns \code{True} --- and there is nothing to leak.

\paragraph{Solution to Exercise \ref{ex:3-reserved}.}
With true signal $R^2 = 0.01$: model (i) scores in- and out-of-sample
$R^2 \approx 0.01$. Model (ii) --- target among its inputs --- scores
$R^2 \approx 1.0$ \emph{both} in-sample and out-of-sample, because the
pipeline that packs the target into the last column does so for test data
too: the leak follows the data, and no holdout detects it. That is the
chilling part: ordinary validation is defenseless against a layout bug.
Appending an innocent feature after the reserved column shifts the slice:
model (i)'s ``correct'' \code{[:, :-1]} now \emph{includes} the target
(instant $R^2 \approx 1$) while dropping a real feature --- one appended
column flipped a correct model into a leaking one without touching its
code. The two-line defense:
\begin{lstlisting}
if x.shape[-1] != len(schema.feature_names):
    raise ValueError(f"width {x.shape[-1]} != schema {len(schema.feature_names)}")
\end{lstlisting}
plus the structural fix of never storing the target among features --- at
which point all three configurations fail loudly at construction, which is
the only acceptable behavior.

\paragraph{Solution to Exercise \ref{ex:3-taxonomy}.}
(a) Leakage channel (near-duplicate rows with overlapping labels cross the
split): defense --- a date-keyed splitting API that offers no random row
split. (b) Leakage channel (test-period moments enter the transform):
defense --- fit-on-train-only as a pipeline contract, pinned by a
shifted-history invariant test. (c) Leakage channel (future positions
visible in training): defense --- a causality test asserting predictions
at $t$ are invariant to perturbations of inputs after $t$ (the package
runs exactly this for the HMM filter). (d) Contract failure (three copies
of one truth): defense --- a single validated config tree where derived
quantities are computed, not restated. (e) Leakage channel (reward peeks
at prices the state already contains): defense --- point-in-time
discipline applied to the \emph{reward} path, same invariant test family
as (b).

\paragraph{Solution to Exercise \ref{ex:3-keep}.}
(Model answers; yours may be better.) \emph{The brief against:} the
skeleton multiplies the parameter and failure-mode budget of a linear
baseline by an order of magnitude --- an encoder, a gate, $K$ experts, a
multimodal likelihood with symmetric local optima and permutation
ambiguity, and (Chapter~\ref{ch:cross}) a target whose predictable
component is $R^2 \sim 10^{-3}$; a convex, deterministic, one-line ridge
regression captures most linear structure at near-zero variance cost,
trains in milliseconds, cannot be mis-optimized, and by
Proposition~\ref{prop:signflip} the mixture's \emph{categorical} advantage
exists only if sign-flipping state-dependence is real and strong --- an
empirical bet, not a theorem, and the burden of proof sits on the
expensive model. \emph{The rebuttal:} the bet is precisely what the system
is built to settle --- cheaply and fairly: capacity-matched baselines
remain permanently in the comparison, planted-truth worlds verify both
that the mixture wins where regimes exist and that the baselines win where
they don't, and the evaluation protocol (purged walk-forward,
multiplicity-corrected) prices the extra failure modes rather than hiding
them. Keeping the skeleton is defensible \emph{because} the referee is
already hired; keeping it on enthusiasm alone would not be. Part~V holds
the whistle.
